\chapter*{Передмова}\label{ch:preface}
\addcontentsline{toc}{chapter}{Передмова}

\paragraph{Про що цей посібник.} Посібник розповідає про фізику гарячої плазми всередині
токамака: як магнітне поле утримує плазму, як встановлюється рівновага, як переносяться
частинки, енергія та імпульс, чим обмежена стійкість і що відбувається на межі плазми зі
стінкою. Наскрізна тема --- \emph{потоки й обертання плазми} та радіальне електричне поле
$\Er$ (розд.~\ref{ch:flows}); саме вони є предметом проєкту «Моделювання потоку плазми в
контурі токамаку», для якого посібник написано.

\paragraph{На чому він ґрунтується.} Текст спирається на корпус із 29 наукових джерел
(статті, звіти лабораторій, дисертації), зібраних у теці \texttt{library/} проєкту, та на
власні розрахунки проєкту. Кожна формула й кожне число з корпусу мають посилання; «с.~$N$» у посиланні --- це сторінка PDF-файлу в \texttt{library/}, а не друкована сторінка видання.
Там, де корпус мовчить, але без стандартної теорії виклад був би неповним (дрейфи
частинок, банановий режим, межа Троєна тощо), теорію наведено в сірих рамках
«Поза корпусом: стандартна теорія». Такі рамки можна перевірити за будь-яким
підручником фізики плазми, але вони \emph{не} підкріплені джерелами корпусу.

\paragraph{Як читати рамки.}
\begin{itemize}
  \item \textbf{Поза корпусом} (сіра) --- стандартна теорія без посилання на корпус.
  \item \textbf{Приклад з корпусу} (синя) --- реальні установки, розряди й числа.
  \item \textbf{Виведення} --- стиснуте виведення з пропущеними рутинними кроками.
  \item \textbf{Практикум} (зелена) --- завдання на коді проєкту: що запустити і що виміряти.
  \item \textbf{Встановлено в проєкті [E\#]} (бірюзова) --- результат, виміряний власним
        кодом проєкту або прочитаний у відкритому джерелі (реєстр~\cite{ProjRegisters}).
  \item \textbf{Гіпотеза [C\#]} (помаранчева) --- \emph{неперевірене} твердження проєкту.
        Для кожної вказано, що її спростує і скільки коштує перевірка. Це відкриті
        дослідницькі задачі для студентів; подавати їх як факт не можна.
  \item \textbf{Спростовано [R\#]} (червона) --- власні твердження проєкту, які не
        витримали перевірки. Їх залишено як повчальні приклади.
\end{itemize}

\paragraph{Дослідницький трек.} Проєкт веде два реєстри: «встановлене» та «гіпотези»
\cite{ProjRegisters}. Правило просте: гіпотеза переходить у встановлене лише після
перевірки, а спростоване не видаляється. Посібник зберігає це правило. Студент, який
перевірить гіпотезу з помаранчевої рамки, робить внесок у проєкт.

\paragraph{Практичні матеріали.} Практикуми використовують чотири частини проєкту:
розв'язувач рівняння Ґреда--Шафранова \texttt{Tokamak-GS-solver}~\cite{ProjGSsolver};
набір даних і стартовий код конкурсу Fusion Equilibrium Challenge~\cite{FEC2026};
тривимірну візуалізацію рівноваги DIII-D \texttt{tokamak-3d-viz}~\cite{ProjViz};
глибокі розбори та проби новизни~\cite{ProjDeepDives}.
Окремий розд.~\ref{ch:jet} на прикладі JET показує анатомію реального токамака: від проєкту
1975~р. до машини, що запрацювала 1983~р., з власною 3D-реконструкцією~\cite{ProjJET}.
Історію гіпотез проєкту (що ми припускали, що перевірили, що спростували) докладно
викладено в окремому супровідному виданні «Путівник гіпотезами проєкту»~\cite{ProjGuide}.

\paragraph{Одиниці та позначення.} Усюди використовується СІ. Температуру подано в
енергетичних одиницях (еВ, кеВ; стала Больцмана поглинута в $T$). Формули з джерел,
що працюють у гаусовій системі, переведено в СІ, про що сказано в тексті. Позначення
єдині для всього посібника (табл.~\ref{tab:notation}) і узгоджені з кодом проєкту та
з форматом EFIT: $\psi$ --- полоїдальний магнітний потік \emph{на радіан}.

\clearpage
\section*{Довідник позначень}\label{sec:notation}
\addcontentsline{toc}{section}{Довідник позначень}

\begin{longtable}{@{}>{\raggedright\arraybackslash}p{3.7cm}>{\raggedright\arraybackslash}p{6.6cm}>{\raggedright\arraybackslash}p{5.6cm}@{}}
\caption{Основні позначення (СІ)}\label{tab:notation}\\
\toprule
Позначення & Величина & Примітка \\
\midrule
\endfirsthead
\toprule
Позначення & Величина & Примітка \\
\midrule
\endhead
\bottomrule
\endfoot
\multicolumn{3}{@{}l}{\emph{Координати}}\\*
$(R,\tor,Z)$ & циліндричні координати & $\tor$ --- тороїдальний кут \\
$\pol$, $\thetastar$ & полоїдальний кут: геометричний; прямих силових ліній & $\dd\thetastar/\dd\tor = 1/q$ уздовж силової лінії \\
$r$, $\rho$ & малий радіус; нормований радіус $\rho=\sqrt{\Psit/\Psi_{\mathrm{t,b}}}$ & \\
$R_0$, $a$, $\varepsilon=a/R_0$, $A=R_0/a$ & великий, малий радіуси; обернене й пряме аспектні відношення & \\
$\kappa$, $\dup$, $\dlo$ & витягнутість; верхня й нижня трикутність & \\
\multicolumn{3}{@{}l}{\emph{Магнітна конфігурація}}\\*
$\psi$ & полоїдальний потік, \si{\weber\per\radian} & $\psiN=(\psi-\psiax)/(\psib-\psiax)$: 0 на осі, 1 на останній замкненій поверхні (LCFS) \\
$\Psit$ & тороїдальний потік (повний), \si{\weber} & $q=\dd\Psit/(2\pi\,\dd\psi)$ \\
$\psit=\Psit/2\pi$ & тороїдальний потік на радіан, \si{\weber\per\radian} & $q=\dd\psit/\dd\psi$; канонічний імпульс силової лінії (розд.~\ref{ch:particles}) \\
$\Bvec=\grad\psi\times\grad\tor+F\grad\tor$ & магнітне поле & $\BR=-\frac1R\frac{\pa\psi}{\pa Z}$, $\BZ=\frac1R\frac{\pa\psi}{\pa R}$, $\Btor=F/R$, $\Bpol=|\grad\psi|/R$ \\
$F(\psi)=R\Btor$ & полоїдальна струмова функція & \\
$\Bzero$ & вакуумне тороїдальне поле на $R_0$ & \\
$\GSop$ & оператор Ґреда--Шафранова & $\GSop\psi=R\frac{\pa}{\pa R}\bigl(\frac1R\frac{\pa\psi}{\pa R}\bigr)+\frac{\pa^2\psi}{\pa Z^2}$ \\
$\Jtor$ & тороїдальна густина струму & $\GSop\psi=-\mu_0R^2p'(\psi)-FF'(\psi)=-\mu_0R\Jtor$ \\
$q$, $\qnf$, $s$ & запас стійкості; $q$ на $\psiN=0{,}95$; магнітний шир $s=\frac rq\frac{\dd q}{\dd r}$ & \\
$\fsa{\cdot}$ & усереднення по магнітній поверхні & \\
$\unitb=\Bvec/B$ & одиничний вектор уздовж поля & індекси $\parallel$, $\perp$ \\
$\Delta_{\mathrm S}$, $\Lambda=\betap+\li/2-1$ & зсув Шафранова; параметр вертикального поля & \\
$n_{\mathrm d}=-\frac{R}{\BZ}\frac{\pa\BZ}{\pa R}$ & індекс спаду поля & \\
$\mathcal{J}$, $V'=\dd V/\dd\psi$ & якобіан координат; похідна об'єму & \\
\multicolumn{3}{@{}l}{\emph{Параметри плазми}}\\*
$n_s$, $T_s$, $e_s$, $m_s$ & густина, температура, заряд, маса сорту $s$ & $s=e,i,z,\alpha,\dots$ \\
$p=\sum_s n_sT_s$ & тиск & \\
$\Zeff$ & ефективний заряд & \\
$\beta$, $\betap$, $\betaN$ & бета плазми; полоїдальна; нормована & \\
$\li$, $\Ip$, $\Vloop$ & внутрішня індуктивність; струм плазми; напруга обходу & \\
$\tauE$, $W$ & час утримання енергії; запасена енергія & \\
$\nGW$ & густина Грінвальда & \\
$Q=P_{\mathrm{fus}}/P_{\mathrm{ext}}$ & підсилення & $Q_s$ з індексом --- тепловий потік \\
$\ITF$, $\Bleg$, $\Tcoil$ & струм TF; поле на нозі TF; натяг котушки (розд.~\ref{ch:jet}) & $\Btor=\mu_0\ITF/(2\pi R)$ \\
$\dripple$, $\eripple$; $\Rv$, $\Lv$ & гофрування; опір та індуктивність камери & $\eripple=2\dripple$; $\rhoe$ --- пит.\ опір \\
\multicolumn{3}{@{}l}{\emph{Потоки й поля}}\\*
$\uvec_s$; $\upar$, $\upol$, $\utor$ & гідродинамічна швидкість; її компоненти & $\vect v$ --- швидкість окремої частинки \\
$\omtor=\utor/R$ & частота тороїдального обертання & \\
$\vE=\ExB/B^2$ & дрейф $\ExB$ & \\
$\Er$, $\Epar$, $\Phi$ & радіальне й паралельне електричне поле; електростатичний потенціал & \\
$\Omc{s}=e_sB/m_s$, $\omega$ & циклотронна частота; частота хвилі & \\
$\omega_E=\dd\Phi/\dd\psi$, $\mathcal{K}$ & частота обертання $\ExB$; неокласичний коефіцієнт $\upol$ & \\
$c_s$, $\mathcal{M}=\upar/c_s$ & швидкість звуку; число Маха & \\
$\rho_s$, $\rho_m$ & ларморів радіус сорту $s$; густина маси & $\rho_s=v_\perp/|\Omc{s}|$; $\rho$ без індексу --- нормований радіус \\
$\nu_s$, $\nustar$, $\lnL$ & частота зіткнень; зіткненнєвість; кулонівський логарифм & \\
\multicolumn{3}{@{}l}{\emph{Транспорт}}\\*
$\Gamma_s$, $Q_s$, $\Pitor$ & потоки частинок, тепла, тороїдального імпульсу & $q$ --- \emph{лише} запас стійкості \\
$D$, $\Vpinch$ & коефіцієнт дифузії; швидкість пінчу & \\
$\chis{s}$, $\chitor$ & температуропровідність; дифузійність імпульсу & $\chi$ завжди з індексом \\ % lint-ok
$\visc$ & тензор в'язкості & \\
$\nu_{\mathrm{cx}}$, $C_T$ & частота перезарядки; коефіцієнт недифузійного потоку імпульсу & \\
$\gamma_E$, $\gamma_T$ & шир $\ExB$; інкремент турбулентності & \\
\multicolumn{3}{@{}l}{\emph{Утікачі й хвилі}}\\*
$\nre$, $\Ec$, $\ED$ & густина утікачів; поле Коннора--Гасті; поле Драйсера & \\
$\gD$, $\tauav$, $\taure$ & темп Драйсера; час лавини; час утримання утікачів & \\
$\tau_{\mathrm s}$, $\Ecr$, $\fSC$ & час сповільнення, критична енергія пучка; частка каналювання потужності α & \\
$\vA$, $\kpar$, $\npar$ & альфвенівська швидкість; паралельні хвильове число й показник заломлення & \\
$\aJ$, $\bJ$ & параметри профілю струму & замість $\alpha,\beta$ у~\cite{Merlo2011} і в коді~\cite{ProjGSsolver} \\
\end{longtable}

\paragraph{Знак $\psi$.} Він залежить від установки та напрямку струму: дані DIII-D
і MAST у корпусі конкурсу~\cite{FEC2026} мають протилежні конвенції знака
(реєстр проєкту, E15~\cite{ProjRegisters}). Тому в
посібнику всі висновки формулюються через $\psiN$, яка від знака не залежить.
